\chapter{Ensemble Methods}
\label{ch:ensembles}

\chapterguide%
  {Decision trees from \chapref{ch:decision-trees} and the bias--variance/model-selection ideas from \chapref{ch:training}. The variance framework is referenced rather than retaught.}
  {compare bagging, forests, boosting, voting, and stacking; validate diversity and complexity gains.}

An \term{ensemble} combines models whose errors differ, allowing mistakes to partly cancel. Random forests and boosted trees are strong tabular candidates, but validation and operational constraints determine suitability.

\section{Why averaging helps}

Averaging reduces variance when errors are not perfectly correlated.

For $B$ models with error variance $\sigma^2$, independent errors give average variance $\sigma^2/B$. With average pairwise correlation $\rho$, it is approximately
\[
  \rho\,\sigma^2+\frac{(1-\rho)}{B}\,\sigma^2 .
\]
Adding models reduces the second term, not the correlation floor. \term{Bagging} creates diversity through independent fits on perturbed data; \term{boosting} adds sequential corrections.

\begin{intuition}
A large committee helps only if members bring different errors. Repeating the same view adds little information.
\end{intuition}

\section{Bagging: bootstrap aggregating}

Bagging fits learners to \term{bootstrap samples}: $n$ draws with replacement. Each sample omits about $37\%$ of rows, since $(1-1/n)^n\approx e^{-1}$. Deep trees are common base learners; combine predictions by voting or averaging.

\begin{mlfigure}
\begin{tikzpicture}[
  font=\small,
  box/.style={draw, rounded corners=2pt, align=center,
              minimum height=2.2em, minimum width=6.4em, inner sep=3pt, fill=gray!10},
  arr/.style={-{Stealth[length=2mm]}, thick}
]
\node[box] (data) {training data\\($n$ observations)};
\node[box, right=1.05cm of data, yshift=1.35cm] (s1) {bootstrap\\sample 1};
\node[box, right=1.05cm of data] (s2) {bootstrap\\sample 2};
\node[box, right=1.05cm of data, yshift=-1.35cm] (s3) {bootstrap\\sample $B$};
\node[box, right=1.05cm of s1] (t1) {tree 1};
\node[box, right=1.05cm of s2] (t2) {tree 2};
\node[box, right=1.05cm of s3] (t3) {tree $B$};
\node[box, right=1.05cm of t2] (agg) {vote /\\average};
\draw[arr] (data) -- (s1);
\draw[arr] (data) -- (s2);
\draw[arr] (data) -- (s3);
\draw[arr] (s1) -- (t1);
\draw[arr] (s2) -- (t2);
\draw[arr] (s3) -- (t3);
\draw[arr] (t1) -- (agg);
\draw[arr] (t2) -- (agg);
\draw[arr] (t3) -- (agg);
\node at ($(s2)!0.5!(s3)$) {$\vdots$};
\node at ($(t2)!0.5!(t3)$) {$\vdots$};
\end{tikzpicture}
\end{mlfigure}

Each row can be predicted by the trees that never saw it; the resulting \term{out-of-bag (OOB)} score is a free validation estimate. Use it as a validation signal for choosing settings, not as a replacement for the final test.

\section{Random forests}
\label{sec:random-forest}

A \term{random forest} adds random feature subsets at each split, reducing correlation between bagged trees. $\sqrt p$ is a common classification starting point; validate feature fractions for regression and other settings.

Practical profile:

\begin{itemize}
  \item \term{Key hyperparameters:} the number of trees $B$ (more trees make the average steadier until the gain flattens out), the number of features tried per split, and tree depth or minimum leaf size.
  \item \term{Little tuning required:} sensible defaults perform well, which makes random forests the standard ``strong model with ten minutes of effort'' on tabular data.
  \item \term{Extras:} OOB error estimates and feature-importance rankings come nearly free; apply the importance caveat from \secref{sec:tree-feature-importance}.
  \item \term{Costs:} the model is large, slower to predict than a single tree, and no longer readable as rules.
\end{itemize}

\subsection{Extra Trees: randomize the split thresholds too}
\label{sec:extra-trees}

\term{Extremely Randomized Trees} (Extra Trees) go one step further than a forest: at each split they draw the candidate thresholds at random instead of searching for the best one, and by default they use the whole training set rather than a bootstrap sample. The extra randomness lowers variance and training time at the cost of a little bias.

Compare Extra Trees with Random Forest under the same folds and metric. Neither is universally better; the useful lesson is that both trade individual-tree strength for diversity across the ensemble.

\section{Boosting: learning from mistakes}

Boosting adds \term{weak learners}, often shallow trees, in sequence to correct the current ensemble. Unlike independent bagging, each fit depends on previous predictions.

\begin{mltable}
\begin{tabular}{@{}p{0.26\linewidth}p{0.31\linewidth}p{0.31\linewidth}@{}}
\toprule
\rowcolor{MLPaleBlue!92}
 & Bagging / random forest & Boosting \\
\midrule
Training & Parallel, independent & Sequential, each model depends on the last \\
Base models & Fully grown, low-bias, high-variance trees & Shallow, high-bias, low-variance trees \\
Main effect & Reduces variance & Reduces bias (and variance via shrinkage) \\
Overfitting risk & Low; robust to noise & Real; needs a learning rate and early stopping \\
Tuning effort & Small & Moderate to large \\
\bottomrule
\end{tabular}
\end{mltable}

\subsection{AdaBoost: reweight the hard cases}
\label{sec:adaboost}

AdaBoost repeatedly increases the weights of misclassified observations:

\begin{enumerate}
  \item Start with equal weights $w_i=1/n$.
  \item Fit a weak learner $h_m$ to the weighted data and compute its weighted error rate $\varepsilon_m$.
  \item Give the learner a say proportional to its quality:
  \[
    \alpha_m=\tfrac{1}{2}\ln\frac{1-\varepsilon_m}{\varepsilon_m}.
  \]
  \item Multiply misclassified weights by $e^{\alpha_m}$ and correct weights by $e^{-\alpha_m}$, then renormalize.
  \item After $M$ rounds, predict with the weighted vote $\operatorname{sign}\!\left(\sum_m\alpha_m h_m(\bm{x})\right)$.
\end{enumerate}

At $\varepsilon_m=0.3$, $\alpha_m\approx0.42$, multiplying misclassified weights by about $1.5$. A chance learner ($\varepsilon_m=0.5$) gets no vote. Repeated emphasis on mislabeled cases makes AdaBoost sensitive to label noise.

\subsection{Gradient boosting: fit the residuals}
\label{sec:gradient-boosting}

Gradient boosting fits $h_m$ to the negative loss gradient at current predictions. For squared loss this is the residual $y_i-F_{m-1}(\bm x_i)$. Update by
\[
  F_m(\bm{x})
  =F_{m-1}(\bm{x})+\nu\,h_m(\bm{x}),
  \qquad 0<\nu\leq 1.
\]
The learning rate $\nu$ scales each correction. Tune it jointly with tree count, using validation or early stopping.

\begin{workedexample}
For $y=(3,5,8,10)$, $F_0=6.5$ gives residuals $(-3.5,-1.5,1.5,3.5)$. A stump grouping the first two and last two rows predicts residual means $-2.5$ and $2.5$. With $\nu=0.5$,
\[
  F_1=(6.5-1.25,\;6.5-1.25,\;6.5+1.25,\;6.5+1.25)
  =(5.25,\,5.25,\,7.75,\,7.75),
\]
and residuals become $(-2.25,-0.25,0.25,2.25)$. Later rounds add further corrections.
\end{workedexample}

\begin{mlfigure}
\begin{tikzpicture}[
  x=1cm, y=1cm, font=\small,
  ax/.style={draw=MLMuted!55,semithick},
  obs/.style={circle,fill=MLNavy,inner sep=1.5pt},
  fit/.style={draw=MLBlue,line width=1.1pt},
  resid/.style={-{Stealth[length=1.7mm]},draw=MLRed,semithick},
  rl/.style={font=\tiny\color{MLRed}},
  hd/.style={font=\small\bfseries\color{MLNavy}},
  yl/.style={font=\tiny\color{MLMuted},anchor=east}]

% ---------- round 0: predict the mean ----------
\node[hd] at (1.9,4.35) {Round 0: predict the mean};
\draw[ax] (0.30,0) -- (0.30,3.95);
\draw[ax] (0.30,0) -- (3.55,0);
\foreach \v/\cy in {3/0.42, 5/1.26, 8/2.52, 10/3.36} \node[yl] at (0.22,\cy) {\v};

\draw[fit] (0.45,1.89) -- (3.35,1.89);
\node[font=\tiny\color{MLBlue},anchor=west] at (3.40,1.89) {$F_0=6.5$};

\draw[resid] (0.70,1.89) -- (0.70,0.50);
\draw[resid] (1.50,1.89) -- (1.50,1.34);
\draw[resid] (2.30,1.89) -- (2.30,2.44);
\draw[resid] (3.10,1.89) -- (3.10,3.28);
\node[rl,anchor=west] at (0.78,1.16) {$-3.5$};
\node[rl,anchor=west] at (1.58,1.62) {$-1.5$};
\node[rl,anchor=west] at (2.38,2.16) {$+1.5$};
\node[rl,anchor=west] at (3.18,2.62) {$+3.5$};

\foreach \cx/\cy in {0.70/0.42, 1.50/1.26, 2.30/2.52, 3.10/3.36}
  \node[obs] at (\cx,\cy) {};

% ---------- round 1: add half of a stump ----------
\begin{scope}[xshift=7.0cm]
\node[hd] at (1.9,4.35) {After round 1, with $\nu=0.5$};
\draw[ax] (0.30,0) -- (0.30,3.95);
\draw[ax] (0.30,0) -- (3.55,0);
\foreach \v/\cy in {3/0.42, 5/1.26, 8/2.52, 10/3.36} \node[yl] at (0.22,\cy) {\v};

\draw[fit] (0.45,1.365) -- (1.90,1.365);
\draw[fit] (1.90,2.415) -- (3.35,2.415);
\draw[fit,dashed] (1.90,1.365) -- (1.90,2.415);
\node[font=\tiny\color{MLBlue},anchor=west] at (3.40,2.415) {$F_1$};

\draw[resid] (0.70,1.365) -- (0.70,0.50);
\draw[resid] (1.50,1.365) -- (1.50,1.34);
\draw[resid] (2.30,2.415) -- (2.30,2.44);
\draw[resid] (3.10,2.415) -- (3.10,3.28);
\node[rl,anchor=west] at (0.78,0.86) {$-2.25$};
\node[rl,anchor=west] at (1.58,1.30) {$-0.25$};
\node[rl,anchor=west] at (2.38,2.05) {$+0.25$};
\node[rl,anchor=west] at (3.18,2.85) {$+2.25$};

\foreach \cx/\cy in {0.70/0.42, 1.50/1.26, 2.30/2.52, 3.10/3.36}
  \node[obs] at (\cx,\cy) {};
\end{scope}

\node[font=\scriptsize\color{MLInk},align=center,text width=12.5cm] at (5.4,-0.85)
  {The stump is never fitted to the labels --- it is fitted to the red arrows.
   One split separates the first two observations from the last two and predicts
   their residual means, $-2.5$ and $+2.5$; the learning rate admits only half
   of that, so every residual shrinks but none is erased. A hundred such
   corrections trace the data closely.};

\end{tikzpicture}
\end{mlfigure}

\begin{intuition}
Gradient boosting fits corrections; AdaBoost reweights mistakes. Both train sequentially.
\end{intuition}

\subsection{Modern gradient-boosting libraries}

Three refinements of gradient boosting dominate practical work; all keep the algorithm above and add engineering and regularization:

\begin{mltable}
\begin{tabularx}{\linewidth}{@{}>{\bfseries\leavevmode\color{MLNavy}}p{0.18\linewidth}p{0.46\linewidth}X@{}}
\toprule
\rowcolor{MLPaleBlue!92}
Library & Distinctive approach & Practical strength \\
\midrule
XGBoost & Adds explicit regularization of tree size and leaf values to the objective and uses second-derivative information for better splits. & Popularized the modern toolkit, including early stopping and column subsampling. \\
LightGBM & Grows trees leaf-wise instead of level-wise and buckets feature values into histograms. & Extremely fast on large datasets. \\
CatBoost & Handles categorical features natively with label statistics computed in a careful order to avoid leakage. & Notably strong when the data contains many categorical columns. \\
\bottomrule
\end{tabularx}
\end{mltable}

The hyperparameters that matter most, in rough order: number of trees (via early stopping), learning rate, tree depth or leaf count, and the row/column subsampling fractions.

\section{Bagging and boosting, side by side}
\label{sec:bagging-vs-boosting}

Bagging fits independently and can run in parallel; boosting fits sequential corrections.

\begin{mlfigure}
\begin{tikzpicture}[
  node distance=0.3cm and 0.55cm,
  t/.style={draw=MLBlue,fill=MLPaleBlue,rounded corners=1.2mm,minimum width=1.45cm,minimum height=0.55cm,align=center,font=\scriptsize\bfseries\color{MLNavy}},
  tb/.style={draw=MLOrange,fill=MLPaleOrange,rounded corners=1.2mm,minimum width=1.45cm,minimum height=0.55cm,align=center,font=\scriptsize\bfseries\color{MLNavy}},
  d/.style={draw=MLMuted,fill=MLPaleGray,rounded corners=1.2mm,minimum width=1.1cm,minimum height=0.5cm,align=center,font=\tiny\color{MLNavy}},
  agg/.style={draw=MLGreen,fill=MLPaleTeal,rounded corners=1.5mm,minimum width=1.5cm,minimum height=1.9cm,align=center,font=\scriptsize\bfseries\color{MLNavy}},
  result/.style={draw=MLGreen,fill=MLPaleTeal,rounded corners=1.2mm,minimum width=1.15cm,minimum height=0.55cm,align=center,font=\scriptsize\bfseries\color{MLNavy}},
  arr/.style={-{Stealth[length=1.8mm]},thick,draw=MLTeal},
  arro/.style={-{Stealth[length=1.8mm]},thick,draw=MLOrange}
]
% ---- bagging ----
\node[font=\small\bfseries\color{MLNavy}] (hb) at (0,2.55) {Bagging: parallel and independent};
\node[d] (d1) at (-1.3,1.85) {sample 1};
\node[d] (d2) at (-1.3,1.15) {sample 2};
\node[d] (d3) at (-1.3,0.45) {sample 3};
\node[t] (t1) at (0.55,1.85) {tree 1};
\node[t] (t2) at (0.55,1.15) {tree 2};
\node[t] (t3) at (0.55,0.45) {tree 3};
\node[agg] (av) at (2.5,1.15) {average\\or vote};
\draw[arr] (d1)--(t1); \draw[arr] (d2)--(t2); \draw[arr] (d3)--(t3);
\draw[arr] (t1)--(av); \draw[arr] (t2)--(av); \draw[arr] (t3)--(av);
\node[font=\scriptsize\itshape\color{MLMuted}] at (0.6,-0.2) {reduces variance};
% ---- boosting ----
\node[font=\small\bfseries\color{MLNavy}] (hs) at (8.85,2.55) {Boosting: sequential and corrective};
\node[tb] (b1) at (5.75,1.45) {tree 1};
\node[tb] (b2) at (8.05,1.45) {tree 2};
\node[tb] (b3) at (10.35,1.45) {tree 3};
\node[result] (sum) at (12.15,1.45) {weighted\\sum};
\draw[arro] (b1.east)--(b2.west);
\draw[arro] (b2.east)--(b3.west);
\draw[arro] (b3.east)--(sum.west);
\node[font=\tiny\color{MLRed},align=center] at (6.90,0.88)
  {fit remaining\\errors};
\node[font=\tiny\color{MLRed},align=center] at (9.20,0.88)
  {fit remaining\\errors};
\node[font=\scriptsize\itshape\color{MLMuted}] at (8.85,-0.2) {reduces bias};
\end{tikzpicture}
\end{mlfigure}

\begin{keyidea}
Bagging often uses deep trees; boosting often uses shallow trees. Validate depth, leaf size, learning rate, and rounds together.
\end{keyidea}

\section{Voting and stacking}

\term{Voting} combines models by majority labels or averaged probabilities. \term{Stacking} trains a \term{meta-model} on base predictions to learn their combination.

Train the meta-model on \term{out-of-fold} base predictions, not predictions fitted to those same rows. Otherwise it learns from overfitted outputs. Linear or logistic meta-models are common.

\section{Practical ensemble workflow}

\begin{itemize}
  \item Want a strong tabular model with minimal fuss? \term{Random forest}.
  \item Chasing the best accuracy and willing to tune? \term{Gradient boosting} with early stopping.
  \item Noisy labels? Prefer forests over AdaBoost-style boosting, or lower the boosting learning rate and depth.
  \item Stack diverse models only when out-of-fold validation justifies the added maintenance.
\end{itemize}

In Lab~3, boosting beats forests on Ames Housing but loses on wine quality. Weigh dataset-specific gains against tuning and runtime.

\begin{keyidea}
Bagging averages varied trees; boosting adds dependent corrections. Judge their gains against latency, interpretability, and maintenance costs.
\end{keyidea}

\begin{modelcard}{Tree Ensembles - Quick Guide}
\begin{tabularx}{\linewidth}{@{}>{\bfseries\leavevmode\color{MLNavy}}p{0.22\linewidth}X X@{}}
\toprule
\rowcolor{MLPaleBlue!92}
 & Random forest & Gradient boosting \cr
\midrule
Core idea & Average many varied trees & Add shallow trees sequentially to repair errors \cr
Usually strong for & Reliable general-purpose tabular prediction & High predictive accuracy on structured/tabular data \cr
Main controls & Number of trees, sampled features, depth, leaf size & Number of trees, learning rate, depth, subsampling \cr
Main risks & Large models, less transparent than one tree & Overfitting through too many rounds, sensitive tuning, slower training \cr
\bottomrule
\end{tabularx}
\end{modelcard}

\subsection{Controls worth tuning}

Random forests are usually forgiving: start with enough trees for stable averages, then control individual-tree complexity with depth, leaf size, and the number of features considered at each split. Gradient boosting is more sensitive because learning rate and number of rounds trade off directly. The table below is a practical starting range, not a universal recipe.

\begin{pitfall}
Early-stopping validation must preserve time or group structure. Supply an appropriately separated validation set or tune rounds with structure-preserving folds.
\end{pitfall}

\begin{mltable}
\begin{tabularx}{\linewidth}{@{}>{\bfseries\leavevmode\color{MLNavy}\ttfamily}p{0.21\linewidth}p{0.16\linewidth}X@{}}
\toprule
\rowcolor{MLPaleBlue!92}
\normalfont\bfseries\leavevmode\color{MLNavy}Setting & \normalfont\bfseries\leavevmode\color{MLNavy}Try & \normalfont\bfseries\leavevmode\color{MLNavy}Effect \\
\midrule
learning\_rate & 0.01--0.1 & Smaller learns more carefully; needs more trees. \\
n\_estimators & 300--2000 & With early stopping, set it generously. \\
max\_depth & 3--8 & Deeper catches interactions but overfits sooner. \\
subsample & 0.6--1.0 & Row sampling; adds randomness, fights overfitting. \\
colsample\_bytree & 0.6--1.0 & Column sampling; same idea, per feature. \\
reg\_lambda & 0--10 & Extra penalty on leaf values. \\
\bottomrule
\end{tabularx}
\end{mltable}

\subsection{When combining model families helps}

\begin{intuition}
Stacking can adapt model weights to different cases, but the extra layer needs sufficient data and out-of-fold training.
\end{intuition}

\begin{pitfall}
Combine individually useful, non-redundant models, such as a linear model and a forest. Keep the ensemble only if validation gains justify its complexity.
\end{pitfall}

\section{Inspecting ensemble predictions}
\label{sec:ensemble-interpretability}

Use permutation importance (\secref{sec:tree-feature-importance}), partial dependence, or SHAP to describe model behavior, not causal effects.

Correlated features can mask permutation importance. Partial dependence averages predictions as one feature varies, potentially creating implausible combinations.

\subsection{Optional: SHAP for local explanations}

\term{SHAP} (SHapley Additive exPlanations) decomposes a prediction relative to a reference output:
\[
  \hat{y}_i = \mathbb{E}[\hat{y}] + \sum_{j=1}^{p}\phi_{ij}.
\]
$\phi_{ij}$ attributes row $i$'s departure from the reference output. Correlated features can redistribute attribution; check stability.

Apply the association-versus-causation distinction from \secref{sec:correlation}: an explanation can reveal what the model used without proving what would happen under an intervention.

\section{Ensemble comparison at a glance}

\begin{mltable}
\begin{tabularx}{\linewidth}{@{}>{\bfseries\leavevmode\color{MLNavy}}p{0.22\linewidth}p{0.25\linewidth}X@{}}
\toprule
\rowcolor{MLPaleBlue!92}
Method & What changes from one tree & Main reason to try it \\
\midrule
Bagging & Bootstrap many trees and average/vote & Reduce variance and make a high-variance learner more stable \\
Random Forest & Bagging plus random feature subsets at splits & Make the trees less correlated, so averaging helps more \\
AdaBoost & Train weak learners sequentially with emphasis on previous errors & Reduce bias by focusing later learners on difficult cases \\
Soft voting & Average probabilities from different model families & Combine complementary decision rules rather than many copies of one rule \\
\bottomrule
\end{tabularx}
\end{mltable}

\begin{tryit}
Compare a single tree, a random forest, and a gradient-boosting model with the same cross-validation splits. Record mean validation score, descriptive fold spread, fit time, and prediction time. Decide whether any gain justifies the added cost, then evaluate the locked choice on the test set once.
\end{tryit}

\begin{quickcheck}
\begin{enumerate}
  \item How does bagging reduce variance, and how does boosting differ in the way later learners are constructed?
  \item Why does a random forest consider only a subset of features at many splits?
  \item Why must predictions used to train a stacking meta-model be out-of-fold rather than fitted predictions from the same rows?
\end{enumerate}
\end{quickcheck}
